\documentclass[uplatex,dvipdfmx,a4paper,11pt]{jsarticle}
\usepackage[margin=20mm]{geometry}
\usepackage{fancyvrb}
\usepackage[dvipdfmx]{hyperref}
\title{プログラミング演習III 課題C1 一括回答レポート}
\date{生成日時: 2026-04-09 13:19}
\begin{document}
\maketitle
\tableofcontents
\newpage
\section*{課題C1 一括回答}
\addcontentsline{toc}{section}{課題C1 一括回答}
\begin{Verbatim}[fontsize=\small,frame=single]
作業ディレクトリ: /home/s24/t246270/C3/1/Prog3/C1
\end{Verbatim}
\section*{例題2: hello.c の出力}
\addcontentsline{toc}{section}{例題2: hello.c の出力}
\begin{Verbatim}[fontsize=\small,frame=single]
Hello World!
\end{Verbatim}
\section*{例題3: repeat.c の出力}
\addcontentsline{toc}{section}{例題3: repeat.c の出力}
\begin{Verbatim}[fontsize=\small,frame=single]
(1) ./myrepeat Hello World!
[0] ./myrepeat [1] Hello [2] World! 
(2) ./myrepeat A B C
[0] ./myrepeat [1] A [2] B [3] C
\end{Verbatim}
\section*{例題4: functions.c の出力}
\addcontentsline{toc}{section}{例題4: functions.c の出力}
\begin{Verbatim}[fontsize=\small,frame=single]
100
1330
\end{Verbatim}
\section*{例題5: substitute.c の出力}
\addcontentsline{toc}{section}{例題5: substitute.c の出力}
\begin{Verbatim}[fontsize=\small,frame=single]
i=1 j=2
i=10 j=100
\end{Verbatim}
\section*{例題6: malloc1.c 空欄と実行}
\addcontentsline{toc}{section}{例題6: malloc1.c 空欄と実行}
\begin{Verbatim}[fontsize=\small,frame=single]
空欄(1): (int *)malloc(sizeof(int) * n)
A[0] = 0
A[1] = 1
A[2] = 4
A[3] = 9
A[4] = 16
A[5] = 25
A[6] = 36
A[7] = 49
A[8] = 64
A[9] = 81
\end{Verbatim}
\section*{例題7: read\_array.c 空欄と実行}
\addcontentsline{toc}{section}{例題7: read\_array.c 空欄と実行}
\begin{Verbatim}[fontsize=\small,frame=single]
空欄(1): read_array(argv[1], n)
空欄(2): (int *)malloc(sizeof(int) * n)
123 456 789 1011 1213 1415 1617 1819 
8443
\end{Verbatim}
\section*{基本問題1: test.c (入力27)}
\addcontentsline{toc}{section}{基本問題1: test.c (入力27)}
\begin{Verbatim}[fontsize=\small,frame=single]
Input an integer: The square is 729.
The cube is 19683.
\end{Verbatim}
\section*{基本問題2: 各コマンドで作られるファイル}
\addcontentsline{toc}{section}{基本問題2: 各コマンドで作られるファイル}
\begin{Verbatim}[fontsize=\small,frame=single]
(1) gcc -c main.c -> main.o
(2) gcc -c interface.c -> interface.o
(3) gcc -c arithmetic.c -> arithmetic.o
(4) gcc main.o interface.o arithmetic.o -o mytest -> mytest
\end{Verbatim}
\section*{基本問題3: make の出力(4行)}
\addcontentsline{toc}{section}{基本問題3: make の出力(4行)}
\begin{Verbatim}[fontsize=\small,frame=single]
gcc -c main.c
gcc -c interface.c
gcc -c arithmetic.c
gcc main.o interface.o arithmetic.o -o mytest
\end{Verbatim}
\section*{基本問題4: 端末に入力する内容}
\addcontentsline{toc}{section}{基本問題4: 端末に入力する内容}
\begin{Verbatim}[fontsize=\small,frame=single]
(1) make
(2) make clean
(3) make interface.o
\end{Verbatim}
\section*{基本問題5: make\_array のmakefile内容と実行}
\addcontentsline{toc}{section}{基本問題5: make\_array のmakefile内容と実行}
\begin{Verbatim}[fontsize=\small,frame=single]
--- makefile ---
OBJS = main.o array_operations.o
CC = gcc

array_test: $(OBJS)
	$(CC) $(OBJS) -o array_test

main.o: main.c C1.h
	$(CC) -c main.c

array_operations.o: array_operations.c C1.h
	$(CC) -c array_operations.c

clean:
	rm -f $(OBJS) array_test
gcc -c main.c
gcc -c array_operations.c
gcc main.o array_operations.o -o array_test
123 456 789 1011 1213 1415 1617 1819
\end{Verbatim}
\section*{発展課題1: 最頻値}
\addcontentsline{toc}{section}{発展課題1: 最頻値}
\begin{Verbatim}[fontsize=\small,frame=single]
サンプル array.txt の最頻値:
3
\end{Verbatim}
\section*{後片付け}
\addcontentsline{toc}{section}{後片付け}
\begin{Verbatim}[fontsize=\small,frame=single]
(出力なし)
\end{Verbatim}
\section*{完了}
\addcontentsline{toc}{section}{完了}
\begin{Verbatim}[fontsize=\small,frame=single]
一括回答の生成が完了しました。
\end{Verbatim}
\end{document}
